% Notes-only two-parameter extension of D.93. Unregistered candidate.
\subsection{Ordinary points with an independently varying local index}

\begin{statement}{D.94}{Open}{Two parameters of finite torsion and free compact kernel}
Let \(N\ge r\ge1\), put \(h=N-r\), and take the actual
rational point of the tripod
\[
 c=5^{h+1}-1,\qquad t_{r,N}=c^{5^r},\qquad
 s_{r,N}=(\kappa(t_{r,N}),\kappa(1-t_{r,N})).
                                                        \tag{D.94.1}
\]
Use the geometrically pro-five metabelian quotient and its actual
Alexander coefficient module of D.35. The entire compact term
\(\mathcal S_{s_{r,N}}\) and separated compact second-cohomology
kernel \(\mathcal K_{s_{r,N}}\) of D.63 have abstract compact
\(\mathbb Z_5\)-module decompositions
\[
 \begin{aligned}
 \mathcal S_{s_{r,N}}
   &\simeq\left(\bigoplus_{j=1}^{r-1}
                     \mathbb Z/5^{h+j}\right)
                       \oplus\mathbb Z/5^{N+1},\\
 \mathcal K_{s_{r,N}}
   &\simeq\mathcal S_{s_{r,N}}\oplus\mathbb Z_5^h.
 \end{aligned}                                         \tag{D.94.2}
\]
The finite sum is empty for \(r=1\). In particular the point
\(t=24^5\) has \(\mathcal S\simeq\mathbb Z/125\) and
\(\mathcal K\simeq\mathbb Z/125\oplus\mathbb Z_5\).
These decompositions do not replace the natural compact-kernel
map: that map is computed explicitly below.

Every indicated metabelian section fibre is nonempty and has
translation group \(\prod_{\mathbb N}\mathbb Z_5\). It contains
exactly one ordinary rational-point section class and no cusp
section class. The assertion gives no higher-derived lift of the
other metabelian classes and no classification for arbitrary
infinitely ramified coefficient pairs. For \(N=r\), its module
calculation is the D.93 calculation, although the displayed
ordinary point may differ from the point used there.
\end{statement}

\begin{proof}[Complete candidate proof]
\paragraph{An actual ordinary point and its exact two lattices.}
Write \(a=5^r\), \(n=5^N\), and \(k=5^h=n/a\).
The integer \(c=5^{h+1}-1\) is not a fifth power. Otherwise
\(c=d^5\), with \(d>1\), would give
\[
 5^{h+1}=(d+1)(d^4-d^3+d^2-d+1).
\]
Reduction modulo five gives \(d\equiv-1\pmod5\). The second
factor has five-adic valuation exactly one, by expanding at
\(d=-1\), but it is strictly larger than five for \(d>1\).
It therefore cannot be a power of five, a contradiction.
Consequently some prime valuation of \(c\) is a five-adic unit.
The integer \(c^a-1\) is not a fifth power either: with
\(d=c^{5^{r-1}}>1\), it lies strictly between \((d-1)^5\)
and \(d^5\). Its prime support is disjoint from that of \(c\).
Neither integer is divisible by five, since \(c\equiv-1\) and
\(c^a-1\equiv-2\pmod5\). Thus the closed away-five valuation
lattice of the actual pair is exactly
\[
             L=5^r\mathbb Z_5\oplus\mathbb Z_5,
                    \qquad (a_5,b_5)=(0,0).            \tag{D.94.3}
\]
The sign in \(1-t\) has zero pro-five Kummer class.

Use the integral local logarithm coordinate
\(\lambda_x=\log(x^4)/(4\log6)\) of D.78. Since
\(c=-(1-5^{h+1})\), one has
\(v_5(\lambda_c)=h\). Moreover
\[
                 1-c^{5^r}\equiv2\pmod{25}.           \tag{D.94.4}
\]
For \(h\ge1\), this follows from \(c\equiv-1\pmod{25}\).
For \(h=0\), \(c=4\) and \(4^{5^r}\equiv-1\pmod{25}\).
The already checked equality \(\lambda_2\equiv2\pmod5\)
shows that the second local coefficient is a unit. Write the
local line vector as
\[
        z=(5^N\alpha,\beta),\qquad
                       \alpha,\beta\in\mathbb Z_5^\times.
                                                        \tag{D.94.5}
\]
The first, second and diagonal global depths are \((r,0,0)\);
the corresponding local depths are \((N,0,0)\). The diagonal
local coefficient is a unit because \(N\ge1\). D.77 therefore
gives precisely the free boundary factor \(\mathbb Z_5^h\)
in (D.94.2). It remains to compute the entire \(\mathcal S\).

\paragraph{The natural compact map reduces to two actual finite modules.}
Let \(\Lambda=\mathbb Z_5[[u,v]]\), \(t_1=1+u\),
\(t_2=1+v\), \(w=t_1t_2-1\), and \(P=uvw\).
By D.87 the actual global and local operator images are
\(L\rtimes I_0\) and \(\mathbb Z_5z\rtimes I_0\), with the
same graph complement. Here that graph has no valuation term,
by (D.94.3). Define
\[
 R=\Lambda/(t_1^{5^N\alpha}t_2^\beta-1),\qquad
 J=(t_1^a-1,t_2-1)R.
                                                        \tag{D.94.6}
\]
Because \(\beta\) is a unit, eliminate \(t_2\):
\[
 R=\mathbb Z_5[[u]],\quad t=t_1,\quad
 t_2=(t^n)^\gamma,\quad \gamma=-\alpha/\beta\in\mathbb Z_5^\times.
                                                        \tag{D.94.7}
\]
Put \(q_j=t^{5^j}-1\). Then \(v=q_N\) times a unit,
\(w=u\) times a unit, and
\[
                 PR=u^2q_NR,\qquad J=q_rR.             \tag{D.94.8}
\]
The latter identity uses \(r\le N\). The exact D.87 formula,
including its actual inclusion map, is
\[
 \mathcal S_s=
 \ker\left((J/PJ)(-1)_\Gamma\longrightarrow
                   (R/PR)(-1)_\Gamma\right),
                 \qquad\Gamma=\mathbb Z_5^\times.       \tag{D.94.9}
\]
The primitive actual cubic operator of D.77 and D.39 is what
identifies the difference ideal with \((P)\); no formal
substitution for that arithmetic input is being made.

Retain \(q_r\) as generator of the source. Its scalar module
is \(R/(u^2q_N)\), with action
\(\chi^{-1}S_\chi(t^a)\rho_0\), where
\(S_\chi(t)=(t^\chi-1)/(t-1)\). The two low augmentation
jets have tame characters \(\chi^2,\chi^3\), so both their
invariants and their coinvariants, and their other
\(\Gamma\)-homology groups, vanish by exact order-four
averaging. On the retained \(u^2\) layer the module is
\[
 \mathscr B_{r,N}=R_N:=\mathbb Z_5[t]/(t^n-1),\qquad
       gF=S_{\chi(g)}(t^a)S_{\chi(g)}(t)^3F(t^{\chi(g)}).
                                                        \tag{D.94.10}
\]
Indeed modulo \(q_N\) one has \(t_2=1\), and the actual
beta factor is exactly one there by D.39. The two factors
from the retained \(u^2\) supply the remaining powers of
\(S_\chi(t)\). This computes the full source coinvariants.

The target has low tame characters \(\chi,\chi^2\).
Its retained \(u^2\) layer has the action
\[
          \mathscr C_N:\quad S_\chi(t)^3F(t^\chi).
                                                        \tag{D.94.11}
\]
The natural map between (D.94.10) and (D.94.11) is
\emph{multiplication by \(q_r\)}. This follows directly
from \(J\hookrightarrow R\), and the intertwiner identity
\(q_rS_\chi(t^a)=\sigma_\chi(q_r)\) proves equivariance.

The target coinvariants are torsion-free of rank \(N\).
To verify this integrally, multiplication by \(u^3\)
identifies its coinvariants with those of the ordinary
permutation module \(I^3\subset R_N\), since its kernel
is the norm line of tame character \(\chi^3\).
The quotient \(I/I^3\) has tame characters \(\chi,\chi^2\),
so \((I^3)_\Gamma=I_\Gamma\). In
\(0\to I\to R_N\to\mathbb Z_5\to0\), the map on
first homology to \(H_1(\Gamma,\mathbb Z_5)\) is onto:
the fixed basis vector \(1\in R_N\) maps to one.
Thus \(I_\Gamma\) is the augmentation kernel of the free
orbit module \((R_N)_\Gamma\), of rank \(N\).

Over a splitting field of characteristic zero, each primitive
root-order orbit \(5^j\), \(1\le j\le N\), gives one
free target coinvariant. In the source its weighted orbit
product is \(2^{4\cdot5^{j-1}}\ne1\) for \(j\le r\),
and one for \(j>r\); the latter assertion follows by
\(\prod(1+\zeta^{a2^i})=1\) over a complete orbit.
The trivial root also contributes no source coinvariant.
The source has free rank \(h\), and on each free orbit
\(j>r\), the actual multiplier \(q_r\) is nonzero.
The map is therefore injective on the rational free quotient.
Since the target is torsion-free, (D.94.9) is exactly
\[
                  \mathcal S_s=\operatorname{Tor}_{\mathbb Z_5}
                              (\mathscr B_{r,N})_\Gamma.
                                                        \tag{D.94.12}
\]
This is an actual-map argument, not a comparison of coranks.

\paragraph{A boundary module whose integral cyclic blocks are known.}
On the same underlying finite free module \(R_N\), put
\[
              \mathscr A_{r,N}:\quad
                       S_\chi(t^a)F(t^\chi).
                                                        \tag{D.94.13}
\]
Multiplication by \(q_r\) identifies it with
\(q_r\mathbb Z_5[[t-1]]/q_Nq_r\mathbb Z_5[[t-1]]\),
with ordinary substitution. This is exactly the infinite
boundary module in D.91.25, with its parameters there
\(r,m,N_{\rm old}\) equal to \(r,N,\infty\). It gives
\[
 (\mathscr A_{r,N})_\Gamma\simeq
        \mathbb Z_5^h\oplus
                       \bigoplus_{j=1}^{r}\mathbb Z/5^{h+j}.
                                                        \tag{D.94.14}
\]
When \(h=0\), the action is \(\chi\sigma\), and the
same assertion is the direct weighted-orbit calculation
in D.93. No inverse limit of abstract isomorphisms is used.

We retain the actual blocks in this calculation. The exponent
basis splits into the block divisible by \(a\), and, for each
\(1\le j\le r\), the block of exponents whose reduction
modulo \(a\) has valuation \(r-j\). These blocks are
preserved by (D.94.13): substitution multiplies by a unit
and \(S_\chi(t^a)\) adds multiples of \(a\). The zero
block contributes the free part; the \(j\)-th nonzero block
contributes the single cyclic factor \(\mathbb Z/5^{h+j}\).
For precision, the stabilizer proof of D.91 identifies that
cyclic factor as the cokernel of the relation
\((\text{unit},5^h)\) in \(\mathbb Z/5^j\oplus\mathbb Z_5\).
This describes the natural map on each block, not only its order.

\paragraph{The norm kernel survives the first-homology boundary.}
Write \(\mathcal N_N=1+t+\cdots+t^{n-1}\). Multiplication
by \(u^3\) gives the exact equivariant sequence
\[
 0\longrightarrow\mathbb Z_5\mathcal N_N(\chi^4)
 \longrightarrow\mathscr B_{r,N}
 \xrightarrow{u^3} I^3\mathscr A_{r,N}
 \longrightarrow0.
                                                        \tag{D.94.15}
\]
Its integral kernel is exactly the norm line, as in D.93.
The quotient \(\mathscr A_{r,N}/I^3\mathscr A_{r,N}\)
has the three tame characters \(\chi,\chi^2,\chi^3\),
including any torsion in its positive jets. It has zero
\(\Gamma\)-homology. Thus the only remaining possible
loss of the norm coinvariant \(\mathbb Z/5\) is the map
\[
 H_1(\Gamma,\mathscr A_{r,N})\longrightarrow
                          \mathbb Z/5[\mathcal N_N].
                                                        \tag{D.94.16}
\]
We now prove that this map is zero integrally.

The compact procyclic resolution after tame averaging
identifies the first group with actual \(\Gamma\)-invariants.
Let \(f\) be such an invariant. Then \(g=q_rf\) is invariant
for ordinary substitution on \(R_N\). Its coefficients
are constant on each valuation orbit. On a nonzero coset
modulo \(a\), all exponents have the same valuation; hence
those coefficients of \(g\) are constant. Their sum is
zero, since multiplication by \(q_r=t^a-1\) is the cyclic
difference on that coset. There are \(k\) terms and
\(\mathbb Z_5\) is torsion-free, so they are all zero.
Consequently the coefficients of \(f\) are constant on
each such coset.

The reduction of \(f\) in \(R_r\) is invariant for
\(\chi\sigma\). This invariant group is zero: two acts
as twice a finite permutation. Thus the sum of the
coefficients of \(f\) on each coset is zero. The constants
on the nonzero cosets vanish. We have proved the literal
support statement
\[
                      f=F(t^a),\qquad F\in R_h.
                                                        \tag{D.94.17}
\]
The action on \(F\), writing \(Q=t^a\), is
\(S_\chi(Q)\sigma_\chi\). For \(h=0\) there are no
invariants. For \(h>0\), its first three tame jets have
characters \(\chi,\chi^2,\chi^3\); therefore
\[
                  F=(Q-1)^3B(Q),\qquad B\in R_h.
                                                        \tag{D.94.18}
\]
For the multiplier two, invariance implies that
\((S_2(Q)^4\sigma_2-1)B\) is an integral multiple of
the norm of \(R_h\). Evaluation at \(Q=1\) gives
\[
                         15B(1)\in5^h\mathbb Z_5.
                                                        \tag{D.94.19}
\]
An integral lift of \(f\) through \(u^3\) is
\(b=S_a(t)^3B(t^a)\). Its augmentation has valuation
at least \(3r+h-1\). The difference
\((\mathscr B_{r,N}(2)-1)b\) is an integral multiple
\(z\mathcal N_N\), and evaluation at one gives
\[
                  z=\frac{15b(1)}{5^N},\qquad
                         v_5(z)\ge2r\ge2.
                                                        \tag{D.94.20}
\]
This is zero in the norm coinvariant of order five.
It proves (D.94.16) is zero, including the whole actual
first-homology group, not only a chosen invariant basis.

The exact homology sequence therefore gives
\[
 0\longrightarrow\mathbb Z/5[\mathcal N_N]
 \longrightarrow(\mathscr B_{r,N})_\Gamma
 \longrightarrow(\mathscr A_{r,N})_\Gamma
 \longrightarrow0.
                                                        \tag{D.94.21}
\]
All these are finitely generated compact \(\mathbb Z_5\)-modules.

\paragraph{The actual norm transfer detects the largest extension.}
Set
\[
           F_{r,N}=q_N/q_r=\sum_{i=0}^{k-1}t^{ai}.
                                                        \tag{D.94.22}
\]
Multiplication by this polynomial gives equivariant injections
from the D.93 modules at depth \(r\) into the present modules:
\[
 \mathscr B_{r,r}\longrightarrow\mathscr B_{r,N},
       \qquad \mathscr A_{r,r}\longrightarrow\mathscr A_{r,N}.
                                                        \tag{D.94.23}
\]
Indeed \(\sigma_\chi(F_{r,N})=F_{r,N}\) in \(R_N\),
and \(S_\chi(t^a)F_{r,N}=\chi F_{r,N}\). These identities
also show that the \(u^3\) sequences commute, and
\[
                    F_{r,N}\mathcal N_r=\mathcal N_N
                                                        \tag{D.94.24}
\]
holds literally. In particular their norm coinvariants
are identified, not multiplied by the index.

We specify the map on the other side of the diagram.
There is an exact sequence
\[
 0\longrightarrow\mathscr A_{r,r}
 \xrightarrow{F_{r,N}}\mathscr A_{r,N}
 \xrightarrow{q_r}J_{r,N}\longrightarrow0,
       \qquad J_{r,N}=\ker(R_N\to R_r),
                                                        \tag{D.94.25}
\]
where the last module has ordinary permutation action.
Restrict to the \(j\)-th nonzero residue block described
after (D.94.14). Its first coinvariant is
\(\mathbb Z/5^j\), by D.93; its middle coinvariant is
\(\mathbb Z/5^{h+j}\), by the retained D.91 block.
The last coinvariant is \(\mathbb Z/5^h\): the upper
and lower permutation blocks are transitive, and their
stabilizer inclusion on first homology is multiplication
by \(5^h\). The exact coinvariant sequence is onto at
its right end. Comparing these three finite orders shows
that the first map is injective, with image exactly
\(5^h\) times the middle cyclic group. This establishes
the needed natural map, block by block. It does not assume
injectivity from a rank calculation.

By D.93 there is a lift \(x\) of the largest quotient
factor at depth \(r\) with
\(5^r x=[\mathcal N_r]\), after a unit rescaling.
Its transfer under (D.94.23) projects to \(5^h\epsilon e_r\),
where \(\epsilon\in\mathbb Z_5^\times\) and \(e_r\)
is a generator of the largest cyclic quotient
\(\mathbb Z/5^N\) in (D.94.14). Choose any lift \(y\)
of \(e_r\). Their difference is a norm class of order
five. Multiplication by \(5^r\) kills that difference,
since \(r\ge1\). By (D.94.24),
\[
                         5^N\epsilon y=[\mathcal N_N].
                                                        \tag{D.94.26}
\]
The right side is nonzero by (D.94.21). The lift has exact
order \(5^{N+1}\). This proves the new largest extension
by transfer of the already proved integral D.93 relation.

\paragraph{All remaining factors and the actual fibre.}
Split the free \(\mathbb Z_5^h\) quotient in (D.94.21)
by choosing lifts of a finite basis; the resulting continuous
\(\mathbb Z_5\)-linear section exists because the source
is a finitely generated compact module. Rescale the largest
lift so that \(5^N y_r=[\mathcal N_N]\). For each
\(j<r\), a lift \(y_j\) of the quotient factor of order
\(5^{h+j}\) satisfies
\(5^{h+j}y_j=c_j[\mathcal N_N]\). Replace it by
\[
                           y_j-c_j5^{r-j}y_r.
                                                        \tag{D.94.27}
\]
These are permissible invertible triangular changes on the
mixed-order quotient, and they produce lifts of exact orders
\(5^{h+j}\). They and \(y_r\) generate the finite torsion
submodule; the product of their orders equals the order
provided by (D.94.21). This proves the finite part in
(D.94.2). Equation (D.94.12) and the free boundary calculation
following (D.94.5) identify the entire two compact terms.

The actual rational point supplies a full section and hence
a metabelian lift. Its two Kummer inputs have finite prime
support, so D.74 gives the stated countable product as the
translation group of its nonempty fibre. The completed
pro-five Kummer map on rational numbers has kernel the sign.
If another ordinary point has the same two coordinates,
then \(t'=\pm t\) and \(1-t'=\pm(1-t)\). For the positive
integer \(t>1\), these equations force \(t'=t\).
The two global Kummer coordinates are linearly independent
over \(\mathbb Q_5\), by their disjoint nonempty supports.
Their pair therefore lies on none of the three cusp lines.
This proves the last assertions at exactly the metabelian
level claimed.

\paragraph{Executed support and remaining scope.}
The finite program \path{notes/de_rank_two_local_index_exact_exploration.py}
computes the literal two-parameter matrix at
\((r,N)=(1,1),(1,2),(2,2),(1,3),(2,3),(3,3)\), with dimensions
up to125. It performs Smith elimination modulo \(5^{10}\),
finds every rationally nonzero pivot below that precision,
and has exactly the rational nullity proved above. Thus those
finite five-primary elementary divisors are exact. They are
support for, and not a substitute for, the all-depth norm
transfer and first-homology proof. Further arithmetic,
cohomological and obstruction claims remain outside the
statement. No new source beyond the exact canonical D.91,
D.93 and their already verified Shapiro/cohomology inputs is
being inferred from a secondary reference.
\end{proof}
